\documentclass[11pt]{article}
\usepackage[T1]{fontenc}
\usepackage{textcomp}
\usepackage[margin=0.75in]{geometry}
\usepackage{amsmath,amssymb,amsthm}
\usepackage{array,booktabs}
\usepackage{hyperref}
\setlength{\parindent}{0pt}
\newtheorem{theorem}{Theorem}
\theoremstyle{definition}
\newtheorem{definition}{Definition}
\title{Flow Proof Artifact: prop-09-fourfold-rectangle-unequal-segments}
\author{Generated by \texttt{flow doc proof}}
\date{Flow Proof Book}
\begin{document}
\maketitle
If a straight line is cut into equal and unequal segments, four times the rectangle by the unequal segments together with the square on the half equals the square on the whole.\\[0.5em]
\textbf{Source.} euclid — Elements, Book II, Proposition 9\\[0.5em]
\bigskip
\noindent{\large \textbf{Derived fact 1.} \textit{If a straight line is cut into equal and unequal segments, four times the rectangle by the unequal segments together with the square on the half equals the square on the whole} \hfill the Euclidean plane \cdot Euclid Book II \cdot Proposition 9: four times the rectangle by unequal parts plus the square on the half equals the square on the whole}\par
\smallskip\noindent\textit{Coordinate: the Euclidean plane \cdot Euclid Book II \cdot Proposition 9: four times the rectangle by unequal parts plus the square on the half equals the square on the whole}\par
\smallskip\noindent\textit{Source: euclid — Elements, Book II, Proposition 9}\par
\smallskip\noindent\textit{Built on: proposition 5: the rectangle by unequal parts plus the square on the midpoint equals the square on the half, for Euclid Book II on the Euclidean plane}\par
\medskip
\noindent\textbf{Goal.} If a straight line is cut into equal and unequal segments, four times the rectangle by the unequal segments together with the square on the half equals the square on the whole\par
\noindent$\displaystyle 4 \cdot rect(\text{unequal parts}) + sq(half) = sq(whole)$\par
\medskip
\noindent
\begin{tabular}{@{} >{\raggedright\arraybackslash}p{0.06\textwidth} >{\raggedright\arraybackslash}p{0.47\textwidth} >{\raggedleft\arraybackslash}p{0.41\textwidth} @{}}
\textbf{\#} & \textbf{Proof} & \textbf{Mathematics} \\
\midrule
\textbf{1.} & We prove that proposition 9: four times the rectangle by unequal parts plus the square on the half equals the square on the whole for Euclid Book II the Euclidean plane. & \\[0.45em]
\textbf{2.} & Let whole = straight\_line\_AB. & \\[0.45em]
\textbf{3.} & Let half = segment\_AM. & \\[0.45em]
\textbf{4.} & Let unequal = segments\_AN\_and\_NB. & \\[0.45em]
\textbf{5.} & From step 2, step 3, and step 4, we can deduce that 4 times rect(unequal) plus sq(half) equals sq(whole). & $\displaystyle 4 \cdot rect(unequal) + sq(half) = sq(whole)$ \\[0.45em]
\textbf{6.} & We invoke the derived fact governing Euclid Book II the Euclidean plane: proposition 5: the rectangle by unequal parts plus the square on the midpoint equals the square on the half, for Euclid Book II on the Euclidean plane. & \\[0.45em]
\textbf{7.} & From step 6, this implies 4 times rect(unequal parts) plus sq(half) equals sq(whole). Hence proven. & $\displaystyle 4 \cdot rect(\text{unequal parts}) + sq(half) = sq(whole)$ \\[0.45em]
\bottomrule
\end{tabular}
\label{thm:Geometry-Euclid Book II-proposition_9_four_times_the_rectangle_by_unequal_parts_plus_the_square_on_the_half_equals_the_square_on_the_whole}
\par\medskip\hrule
\end{document}
